\section{Introduction}
\label{sec:intro}

A language model can understand a sentence and still annotate it incorrectly. Annotation requires judgments that follow a project's categories, boundaries, and conventions for ambiguous cases. Large language models (LLMs) offer a promising starting point \citep{gilardi2023chatgpt,he2024annollm}, but their general competence does not ensure adherence to these requirements.

Human annotators face the same problem. Guidelines turn conceptual definitions into decisions that annotators can reproduce \citep{pustejovsky2012natural,halterman2026what}. For complex ontologies, identifying structures, assigning attributes, and connecting them are interdependent decisions. Applying the guideline can also require external resources.

Recent work incorporates guidelines into prompts and refines them against annotation errors \citep{kim2026refining,kim2026diziner}. Yet specifying the required annotation does not necessarily provide an effective procedure for producing it. Long instructions strain guideline following \citep{liu2024lost,panavas2026handbook}, and rewriting them leaves the workflow unchanged. We therefore ask whether guidelines should become an executable \emph{harness}: step instructions, a workflow, and callable tools.
\begin{quote}
\noindent\textbf{H1:} Organizing a guideline into a harness can improve complex annotation beyond direct guideline prompting.
\end{quote}

A stronger annotator may follow instructions better, but its procedure also shapes its performance. This suggests a second route to better annotation:
\begin{quote}
\noindent\textbf{H2:} Refining the harness can yield larger gains than replacing the annotating model.
\end{quote}

In human annotation projects, \emph{moderation} uses a few annotated examples to reveal disagreements and clarify the guideline \citep{pustejovsky2012natural,kim2026refining}. We extend moderation to the harness: structured error feedback lets a builder revise instructions, reorganize steps, and write tools or rules. We ask whether these changes capture reusable decisions:
\begin{quote}
\noindent\textbf{H3:} Limited annotation supervision can produce harness improvements that transfer to unseen texts and other annotators.
\end{quote}

We investigate these hypotheses with \sys, which builds a seed harness from a guideline and recursively improves it through annotation, fine-grained evaluation, and reflection (Figure~\ref{fig:overview}). Our contributions are threefold:
\begin{enumerate}[leftmargin=*,itemsep=1pt,topsep=2pt]
  \item A framework that makes the annotation procedure itself learnable: instructions, steps, and tools are editable through recursive self-improvement, without fine-tuning (\S\ref{sec:method}).
  \item Evidence that improving the harness can matter more than changing the model: gains transfer across annotators, and the learned procedure can surpass task-specific fine-tuned systems under matched evaluation (\S\ref{sec:main}, \S\ref{sec:general}).
  \item An analysis of what self-improvement learns and where it stops: the builder turns recurring conventions into instructions and code and invents resource lookup tools, while difficult semantic judgments and inconsistent annotation practices remain obstacles (\S\ref{sec:how}).
\end{enumerate}
